\documentclass[aspectratio=169]{beamer}

\usetheme{Madrid}
\usecolortheme{default}
\setbeamertemplate{navigation symbols}{}

\usepackage{amsmath}
\usepackage{amssymb}
\usepackage{booktabs}

\title{Trading Strategy Design}
\subtitle{Mean-Variance Allocation and Momentum Signal}
\author{}
\date{}

\begin{document}

\begin{frame}
  \titlepage
\end{frame}

\begin{frame}{1. Strategy Overview}
  \begin{itemize}
    \item The strategy combines two signals:
    \begin{itemize}
      \item Mean-variance allocation signal
      \item Momentum signal
    \end{itemize}
    \item Mean-variance component:
    \begin{itemize}
      \item Estimates expected return and risk
      \item Allocates capital across assets
    \end{itemize}
    \item Momentum component:
    \begin{itemize}
      \item Uses recent price movement and volume information
      \item Ranks assets cross-sectionally
    \end{itemize}
    \item Final strategy blends the two components dynamically.
  \end{itemize}
\end{frame}

\begin{frame}{2. Mean-Variance Strategy}
  \begin{itemize}
    \item Shrinkage-based Markowitz-style allocation model.
    \item Uses recent asset returns to estimate:
    \begin{itemize}
      \item Expected return vector
      \item Covariance matrix
    \end{itemize}
    \item Objective:
    \begin{itemize}
      \item Balance expected return and portfolio variance
      \item Produce risk-aware portfolio weights
    \end{itemize}
    \item Additional portfolio control:
    \begin{itemize}
      \item Single-asset position cap
    \end{itemize}
  \end{itemize}
\end{frame}

\begin{frame}{2. Mean-Variance Formula}
  \[
    \max_{w}
    \quad
    \mu_{\text{shrunk}}^{\top} w
    -
    \frac{\gamma}{2}
    w^{\top}\Sigma_{\text{shrunk}}w
  \]

  \[
    \text{subject to}
    \quad
    w_i \leq \text{position cap}
  \]

  \vspace{0.4cm}

  Required inputs:
  \begin{itemize}
    \item $\mu_{\text{shrunk}}$: shrinkage expected return vector
    \item $\Sigma_{\text{shrunk}}$: shrinkage covariance matrix
    \item $\gamma$: risk-aversion coefficient
    \item Position cap: maximum weight for a single asset
  \end{itemize}
\end{frame}

\begin{frame}{2. Shrinkage Estimation}
  \begin{itemize}
    \item Expected return shrinkage:
  \end{itemize}

  \[
    \mu_{\text{shrunk}}
    =
    (1-\lambda_{\mu})\mu_{\text{sample}}
    +
    \lambda_{\mu}\mu_{\text{target}}
  \]

  \[
    \mu_{\text{target}} = 0
  \]

  \begin{itemize}
    \item Covariance shrinkage:
  \end{itemize}

  \[
    \Sigma_{\text{shrunk}}
    =
    (1-\lambda_{\Sigma})\Sigma_{\text{sample}}
    +
    \lambda_{\Sigma}\Sigma_{\text{diag}}
  \]

  \begin{itemize}
    \item Purpose: reduce noise in short-window estimates.
  \end{itemize}
\end{frame}

\begin{frame}{3. Momentum Strategy}
  \begin{itemize}
    \item Cross-sectional ranking signal.
    \item Based on:
    \begin{itemize}
      \item Short-term price movement
      \item Volume confirmation
    \end{itemize}
    \item Core process:
    \begin{enumerate}
      \item Compute short-horizon price returns.
      \item Standardize returns across assets using z-scores.
      \item Adjust signal by relative volume.
      \item Rank assets and select strongest candidates.
      \item Assign score-based weights with a position cap.
    \end{enumerate}
  \end{itemize}
\end{frame}

\begin{frame}{3. Momentum Features}
  Current feature list:

  \vspace{0.2cm}

  \begin{itemize}
    \item Short-horizon price return
    \item Multi-horizon momentum score
    \item Cross-sectional momentum z-score
    \item Relative volume
    \item Volume-adjusted momentum score
    \item Cross-sectional ranking score
  \end{itemize}
\end{frame}

\begin{frame}{3. Momentum Signal Construction}
  Price momentum:
  \[
    r_{i,h}
    =
    \frac{P_{i,t}}{P_{i,t-h}} - 1
  \]

  Cross-sectional z-score:
  \[
    z_{i,h}
    =
    \frac{r_{i,h} - \bar{r}_{h}}{\sigma(r_h)}
  \]

  Multi-horizon momentum score:
  \[
    m_i
    =
    \frac{1}{H}\sum_{h \in H} z_{i,h}
  \]
\end{frame}

\begin{frame}{3. Volume Adjustment and Ranking}
  Relative volume:
  \[
    RV_i
    =
    \frac{V_{i,t}}{\overline{V}_{i,t}}
  \]

  Volume-adjusted momentum score:
  \[
    s_i
    =
    m_i \cdot f(RV_i)
  \]

  \vspace{0.3cm}

  Portfolio construction:
  \begin{itemize}
    \item Rank assets by final score $s_i$.
    \item Select strongest candidates.
    \item Assign weights according to score strength.
    \item Apply maximum weight limit for each selected asset.
  \end{itemize}
\end{frame}

\begin{frame}{4. Combined Portfolio Construction}
  Both sub-strategies output target weights:

  \[
    w_{\text{mv}}
    \quad \text{and} \quad
    w_{\text{mom}}
  \]

  Final blended portfolio:

  \[
    w_{\text{final}}
    =
    \alpha w_{\text{mom}}
    +
    (1-\alpha)w_{\text{mv}}
  \]

  \vspace{0.3cm}

  Interpretation:
  \begin{itemize}
    \item Larger $\alpha$: stronger momentum influence.
    \item Smaller $\alpha$: stronger mean-variance influence.
    \item Blending, not model switching.
  \end{itemize}
\end{frame}

\begin{frame}{5. Time-Varying Blend}
  \begin{itemize}
    \item $\alpha$ changes over time.
    \item Early stage:
    \begin{itemize}
      \item Mean-variance model has limited historical data.
      \item Momentum signal receives more weight.
    \end{itemize}
    \item Later stage:
    \begin{itemize}
      \item More data is available for return and risk estimation.
      \item Mean-variance allocation receives more weight.
    \end{itemize}
  \end{itemize}

  \vspace{0.2cm}

  \[
    \text{Fast signal-driven model}
    \quad \longrightarrow \quad
    \text{Risk-aware allocation model}
  \]
\end{frame}

\begin{frame}{6. Future Improvements}
  Improvement 1: momentum feature engineering
  \begin{itemize}
    \item Add more predictive features:
    \begin{itemize}
      \item Volatility
      \item Turnover
      \item Short-term reversal
      \item Liquidity
      \item Price acceleration
    \end{itemize}
    \item Use machine learning models to combine features.
    \item Predict future returns or ranking scores.
  \end{itemize}
\end{frame}

\begin{frame}{6. Future Improvements}
  Improvement 2: better strategy integration
  \begin{itemize}
    \item Current method directly blends two complete weight vectors.
    \item Alternative structure:
    \begin{enumerate}
      \item Use momentum model to select a fixed number of candidate stocks.
      \item Run mean-variance optimization only within the selected universe.
      \item Let mean-variance decide final individual weights.
    \end{enumerate}
    \item Interpretation:
    \begin{itemize}
      \item Momentum decides what to trade.
      \item Mean-variance decides how much to allocate.
    \end{itemize}
  \end{itemize}
\end{frame}

\end{document}
